\section{Discussion}
\label{sec:discussion}

\para{{\bf Does merge order matter?} Yes, but through compression.}
Our results confirm that merge sequence matters for language modeling.
It matters through segmentation quality, not because a model is sensitive to merge order as such.
Every \bpb gap in our study follows how many bytes each token carries.
The same-vocabulary controls make this point most clearly.
\randorder and \revorder have exactly \freq's merges, so they \emph{could} produce the same long tokens.
Their different priority prevents this, leaving 59--79\% of the vocabulary unused.
A merge list is therefore not just a set of tokens: the order is part of the tokenizer, and changing it can quietly turn an 8k vocabulary into a 2k one.
Anyone who edits a merge list, for example by inserting domain merges at the top, should check vocabulary utilization and compression afterward.

\para{{\bf Why is fine-tuning robust?}}
With 20k labeled examples, a 23M-parameter model fine-tuned from any of our language models reaches the same accuracy within about 1\%.
We see two likely reasons.
First, short-text classification depends on lexical and semantic cues that the model can recover from any reasonable segmentation, even a nearly character-level one.
Second, plentiful labeled data lets fine-tuning make up for weaker pre-training.
This is consistent with \citet{goldman2024unpacking}, who find that compression matters more for generation than for classification.
It also suggests that work like \citet{schmidt2024tokenization}, which finds downstream differences between tokenizers with similar compression, may need generative or low-data tasks to see large effects.

\para{{\bf Why is domain weighting so weak?}}
Up-weighting PubMed counts by 8 kept 99\% of the top-1k biomedical words whole, up from 92\%, and improved biomedical compression by 7\%.
This bought only 0.8\% biomedical \bpb, at a cost of 0.4\% general \bpb, and no classification gain.
The reason is that the standard tokenizer corpus is already 25\% biomedical, so \freq already covers most biomedical vocabulary.
The hypothesized 5--12\% domain advantage would likely require a much more under-represented domain or a much smaller vocabulary.

\para{Practical implications.}
For small models fine-tuned on classification with plenty of labels, merge order is close to irrelevant, so practitioners can re-weight or reorder merges without measurable cost.
For language modeling or generation, merge order matters, and the rule is simple: prefer the order that keeps common words whole and gives the best bytes per token on the target text, even when the vocabulary is the same.

\subsection{Limitations}
\label{sec:limitations}

\para{Scale.}
Our models have 23M parameters, are trained on 200\,MB of text (51--195M tokens), and use a single vocabulary size (8,192).
Larger models, longer training, or larger vocabularies may change the size of the effects.
In particular, we do not know whether the \bpb gaps grow or shrink with scale.

\para{Equal-text design.}
Our conditions differ in tokens and compute per byte.
We hold text fixed, which is the natural ``same corpus'' comparison.
An equal-compute or equal-byte-context design would separate these factors further.
Because the poor compressors received more compute, we expect such a design to widen, not close, the \bpb gaps.

\para{Tasks.}
Our three tasks are short-text classification with plentiful labels, which likely makes them insensitive to tokenization.
Generation, few-shot prompting, spelling or arithmetic tasks, and low-resource fine-tuning may reveal effects that our tasks hide.

\para{Seeds and domains.}
We use 2 pre-training seeds per tokenizer (1 for the \textsc{-t1} variants), one tokenizer corpus, and one target domain (biomedical).

\para{Definitions of order.}
Our definitions of ``random'' and ``reverse'' are design choices: a count threshold of 50 for the learning rules, and dependency-respecting reorderings for application order.
Other definitions, such as sampling pairs in proportion to frequency, could give milder degradation.

\para{Fixed fine-tuning hyper-parameters.}
We use the same fine-tuning hyper-parameters for every tokenizer.
This is fair, but it could slightly disadvantage tokenizers that produce longer sequences, for example through truncation at 256 tokens.

\para{Broader impact.}
This study uses public datasets and small models, and it trains no system intended for deployment.
Our findings may help practitioners avoid tokenizer edits that silently reduce compression and waste vocabulary, which in turn wastes compute.
We do not foresee direct negative uses.
